% ID: SEL-D0040-OA-E02
% ESTADO: validada. S041, cotejada con examen y criterios.
\textbf{a)} Dibujamos una semirrecta horizontal hasta $0$,
el arco $1-x^2$ entre 0 y 4 y la rama $(x-4)^2+1$ desde 4.
En $4$, el punto $(4,-15)$ queda abierto y $(4,1)$ está incluido.
\begin{center}
\setlength{\unitlength}{8mm}
\begin{picture}(12,8)(-3,-4.5)
\qbezier(-2.5,0)(2.5,0)(7.5,0)\put(7.7,0){$x$}
\qbezier(0,-4)(0,-0.5)(0,3)\put(0,3.1){$f$}
\qbezier(-2.5,0.25)(-1.25,0.25)(0,0.25)
\qbezier(0,0.25)(2,0.25)(3.97,-3.690225)
\qbezier(4,0.25)(5.5,0.25)(7,2.5)
\put(4,0.25){\makebox(0,0){$\bullet$}}
\put(4,-3.75){\makebox(0,0){$\circ$}}
\put(4.2,0.45){$(4,1)$}
\put(4.2,-3.75){$(4,-15)$}
\end{picture}
\end{center}
\textbf{b)} En $0$, los límites y el valor coinciden en 1,
y las pendientes laterales son ambas cero. En $4$, el límite
izquierdo es $-15$ y el derecho vale 1. Hay una discontinuidad
de salto y no existe derivada allí. En los interiores de los
tramos es continua y derivable por ser polinómica. Por tanto
es continua y derivable en $\mathbb R\setminus\{4\}$.

\textbf{c)} En el tramo central $f'(x)=-2x$ y en el derecho
$f'(x)=2x-8$. Resulta $\boxed{f'(1)=-2,\ f'(5)=2}$.
